\bookchapter{บทที่ 62: การพัฒนาอินเทอร์เฟซ — เริ่มต้นกับ SwiftUI}

\chapterintro{อินเทอร์เฟซคือทุกสิ่งที่ผู้ใช้เห็น SwiftUI คือเฟรมเวิร์กอินเทอร์เฟซแบบประกาศ (declarative) ที่ Apple เปิดตัวในปี 2019 — คุณแค่'ประกาศ'ว่าอินเทอร์เฟซหน้าตาเป็นอย่างไร ระบบจะจัดการเรนเดอร์ ปรับให้เข้ากับโหมดมืด และรองรับหน้าจอทุกขนาดให้อัตโนมัติ ถึงปี 2026 SwiftUI กลายเป็นกระแสหลักของโปรเจกต์ใหม่ไปแล้ว ในบทนี้ เราจะมาเรียนรู้การ'วาดภาพ'กัน}

\booksection{แนวคิดของ SwiftUI: แบบประกาศ ไม่ใช่แบบสั่ง}

การพัฒนาอินเทอร์เฟซแบบเก่าเป็นแบบ'สั่งการ': บอกระบบทีละขั้นว่า'สร้างปุ่ม วางที่ (20,30) ตั้งค่าสี' SwiftUI เป็นแบบ'ประกาศ': คุณบรรยายตรง ๆ ว่า'บนอินเทอร์เฟซมีปุ่มสีแดงเขียนว่า “แตะฉัน”' ส่วนจะวาดอย่างไร เมื่อไหร่วาดใหม่ ปล่อยให้ระบบจัดการเอง

\begin{lstlisting}[language=Swift]
import SwiftUI

struct ContentView: View {
    var body: some View {      // body: บอกว่า "หน้าจอหน้าตาเป็นยังไง"
        VStack {               // เรียงตามแนวตั้ง
            Text("สวัสดีชาวโลก!")
                .font(.largeTitle)      // มอดิฟายเออร์: หัวเรื่องใหญ่
            Button("แตะฉัน") {
                print("ถูกแตะแล้ว")
            }
            .buttonStyle(.borderedProminent)  // ปุ่มทึบสีน้ำเงิน
        }
        .padding()
    }
}
\end{lstlisting}

จุดสำคัญ: วิว = struct ที่คอนฟอร์มกับโปรโตคอล View, พร็อพเพอร์ตี้ body ใช้บรรยายอินเทอร์เฟซ, และ .modifier() ตกแต่งแบบต่อทอดกันเป็นโซ่ นี่คือพื้นฐานทั้งหมดของ SwiftUI

\booksection{วิวพื้นฐาน: ข้อความ ปุ่ม ช่องกรอก รูปภาพ}

\begin{center}
\begin{tabularx}{\linewidth}{|X|X|X|}
\hline
\textbf{วิว} & \textbf{ทำหน้าที่อะไร} & \textbf{ตัวอย่าง} \tabularnewline
\hline
Text & แสดงข้อความ & Text("สวัสดี") \tabularnewline
Button & ปุ่ม & Button("แตะฉัน") \{ การทำงาน \} \tabularnewline
TextField & ช่องกรอกข้อความ & TextField("ข้อความนำ", text: \$ข้อความ) \tabularnewline
Image & แสดงรูปภาพ & \texttt{Image(systemName: "star.\allowbreak fill")} \tabularnewline
\hline
\end{tabularx}
\end{center}

\begin{center}\small\color{gray}วิวพื้นฐานทั้ง 4 ของ SwiftUI\end{center}

สังเกต Image(systemName:) — Apple มีคลังไอคอน SF Symbols ในตัว (ไอคอนฟรีหลายพันตัว) แค่พิมพ์ชื่ออย่าง “star.fill” ก็แสดงดาวทึบได้เลย นี่คือของดีที่ซ่อนอยู่ของการพัฒนา iOS: ทำอินเทอร์เฟซแทบไม่ต้องหาไฟล์ไอคอนเลย

\booksection{ชุดจัดวางสามพี่น้อง: VStack, HStack, ZStack}

\begin{itemize}
  \item VStack: เรียงตามแนวตั้ง (เหมือนต่อแถวรับอาหาร)
  \item HStack: เรียงตามแนวนอน (เหมือนนั่งเคียงกัน)
  \item ZStack: วางซ้อนกัน (เหมือนไพ่กองทับกัน เหมาะทำป้ายมุมหรือเลเยอร์บัง)
\end{itemize}

\begin{lstlisting}[language=Swift]
VStack(alignment: .leading, spacing: 12) {
    HStack {
        Text("บิล")
        Spacer()                 // ช่องยืดหยุ่น: ดันฝั่งขวาไปสุด
        Text("-25.5 หยวน").foregroundColor(.red)
    }
    Divider()                    // เส้นแบ่ง
    Text("โน้ต: มื้อกลางวัน")
}
.padding(16)                     // เว้นรอบด้าน
.background(Color(.systemGray6)) // พื้นเทาอ่อน
.cornerRadius(12)                // มุมโค้ง

// ลำดับมอดิฟายเออร์มีผล: padding ก่อน แล้ว background พื้นหลังไม่รวมช่องว่าง
\end{lstlisting}

\begin{tcbwarn}{ข้อควรระวัง: ลำดับของ modifier มีผล}
padding(16).background(สีเทา) กับ background(สีเทา).padding(16) ให้ผลไม่เหมือนกัน — แบบแรกคือ'เว้นระยะก่อนแล้วระบายสี' (พื้นหลังไม่รวมระยะขอบ) ส่วนแบบหลังคือ'ระบายสีก่อนแล้วเว้นระยะ' (พื้นหลังรวมระยะขอบด้วย) ปัญหา'ขอบขาวปริศนา'ที่มือใหม่เจอบ่อยที่สุดก็มาจากจุดนี้ ไม่แน่ใจเมื่อไหร่ให้สลับลำดับดู พรีวิวจะบอกคำตอบทันที

\end{tcbwarn}

\booksection{สเตตและการไบนด์: ทำให้อินเทอร์เฟซ'มีชีวิต'}

อินเทอร์เฟซนิ่ง ๆ ไม่สนุกหรอก กฎของ SwiftUI คือ: อินเทอร์เฟซสะท้อนสเตตอย่างซื่อสัตย์เสมอ พอสเตตเปลี่ยน อินเทอร์เฟซจะวาดใหม่เองอัตโนมัติ

\begin{lstlisting}[language=Swift]
struct CounterView: View {
    @State private var count = 0   // @State: ค่านี้เปลี่ยนเมื่อไร วาดใหม่เมื่อนั้น

    var body: some View {
        VStack {
            Text("คุณแตะไป \(count) ครั้ง")
            Button("แตะฉัน +1") {
                count += 1          // แก้สเตต หน้าจอเด้งเอง
            }
        }
    }
}
\end{lstlisting}

ช่องกรอกข้อความต้องใช้'การไบนด์สองทาง' — สัญลักษณ์ \$ ส่งทั้ง'ค่าและสิทธิ์แก้ไข'เข้าไปพร้อมกัน:

\begin{lstlisting}[language=Swift]
struct InputView: View {
    @State private var text = ""

    var body: some View {
        TextField("พิมพ์อะไรก็ได้", text: $text)  // $ ไบนดิ้งสองทาง
            .textFieldStyle(.roundedBorder)
        Text("คุณพิมพ์ว่า: \(text)")   // อัปเดตสด ๆ ระหว่างพิมพ์
    }
}
\end{lstlisting}

ผู้ใช้ดำเนินการ

สเตตเปลี่ยน (count += 1)

SwiftUI ตรวจจับการเปลี่ยนแปลง

วาด UI ใหม่อัตโนมัติ (body คำนวณใหม่)

\begin{center}\small\color{gray}ภาพที่ 3: การไหลของข้อมูลทางเดียว'สเตต อินเทอร์เฟซ'ของ SwiftUI\end{center}

\booksection{พรีวิว ธีม และแอนิเมชัน: 'ไอเทมเสริมพลัง'ของ SwiftUI}

\begin{itemize}
  \item มาโคร \#Preview: เพิ่ม \#Preview \{ ContentView() \} ไว้ท้ายไฟล์ แคนวาสด้านขวาจะเรนเดอร์แบบเรียลไทม์ — แก้โค้ดปุ๊บเห็นผลปั๊บ ตาเห็นอะไรก็ได้อย่างนั้น
  \item โหมดมืดได้ฟรี: ใช้สีของระบบ (Color(.systemGray6), .primary) แทนการกำหนดสีตายตัว โหมดมืดจะปรับให้เองอัตโนมัติ ไม่ต้องเขียนเพิ่มแม้แต่บรรทัดเดียว
  \item แอนิเมชันจบในบรรทัดเดียว: ใช้ withAnimation \{ count += 1 \} หรือใส่ .animation(.default, value: count) ให้วิว — ทรานซิชันลื่นไหลสุด ๆ
\end{itemize}

\begin{lstlisting}[language=Swift]
struct AnimatedHeartView: View {
    @State private var liked = false

    var body: some View {
        Button {
            withAnimation(.spring(duration: 0.4)) {
                liked.toggle()
            }
        } label: {
            Image(systemName: liked ? "heart.fill" : "heart")
                .font(.system(size: 48))
                .foregroundColor(liked ? .red : .gray)
                .scaleEffect(liked ? 1.3 : 1.0)
        }
    }
}
\end{lstlisting}

\begin{tcbinfo}{เคล็ดลับ: ฟังก์ชันมาก่อน แอนิเมชันมาทีหลัง}
โรคที่มือใหม่เป็นกันบ่อยคือ'อินเทอร์เฟซยังทำไม่ถูกเลย แต่ไปเป็นกังวลกับแอนิเมชันเสียก่อน' ให้ฟังก์ชันรันผ่านก่อน แล้วค่อยใส่ลูกเล่นทั้งหมดตอนท้าย — แอนิเมชันคือของตกแต่งให้งานที่ดีอยู่แล้วเด่นขึ้น แต่ฟังก์ชันคือสิ่งที่ช่วยชีวิตตอนยากจนที่สุด นี่ก็คือหลัก'เรียบง่ายเหนือสิ่งอื่นใด'ที่เรายึดมาตลอด

\end{tcbinfo}

\begin{tcbpractice}{ลงมือทำ: สร้าง'ปุ่มอารมณ์'}
\booksection{สรุป modifier ที่ใช้บ่อย: มีหน้าเดียวในมือ จัดอินเทอร์เฟซได้ไม่ยาก}

\begin{center}
\begin{tabularx}{\linewidth}{|X|X|}
\hline
\textbf{อยากทำอะไร} & \textbf{Modifier} \tabularnewline
\hline
เว้นระยะขอบ & \texttt{.\allowbreak padding(16)} \tabularnewline
มุมโค้ง & \texttt{.\allowbreak cornerRadius(12)} \tabularnewline
สีพื้นหลัง & \texttt{.\allowbreak background(Color(.\allowbreak systemGray6))} \tabularnewline
ฟอนต์และขนาดตัวอักษร & \texttt{.\allowbreak font(.\allowbreak title.\allowbreak bold())}/\texttt{.\allowbreak font(.\allowbreak system(size: 20))} \tabularnewline
สีพื้นหน้า (ข้อความ/ไอคอน) & \texttt{.\allowbreak foregroundColor(.\allowbreak red)} \tabularnewline
กว้างเต็มจอ & \texttt{.\allowbreak frame(maxWidth: .\allowbreak infinity)} \tabularnewline
ขนาดคงที่ & \texttt{.\allowbreak frame(width: 100, height: 44)} \tabularnewline
การจัดตำแหน่ง & \texttt{VStack(alignment: .\allowbreak leading)} \tabularnewline
ความโปร่งใส & \texttt{.\allowbreak opacity(0.\allowbreak 6)} \tabularnewline
เงา & \texttt{.\allowbreak shadow(radius: 5)} \tabularnewline
เจสเชอร์การแตะ & \texttt{.\allowbreak onTapGesture \{ .\allowbreak .\allowbreak .\allowbreak  \}} \tabularnewline
ชนิดคีย์บอร์ด & \texttt{.\allowbreak keyboardType(.\allowbreak decimalPad)} \tabularnewline
\hline
\end{tabularx}
\end{center}

\begin{center}\small\color{gray}สรุป modifier ที่ใช้บ่อยของ SwiftUI\end{center}

จำไม่ได้ก็ไม่เป็นไร — หน้านี้ไว้เปิดค้นเฉย ๆ จริง ๆ ที่ใช้บ่อยที่สุดมีแค่สี่ตัวแรก: padding, มุมโค้ง, สีพื้นหลัง, ฟอนต์

ลองแปลงโปรเจกต์เทมเพลตเป็นตัวสลับอารมณ์: มีปุ่มหนึ่งปุ่ม แตะแล้ววนสลับไปมาระหว่าง ดีใจ เศร้า โมโห โดยข้อความและสีพื้นหลังเปลี่ยนตาม คำใบ้: ใช้อาร์เรย์เก็บสามชุด (ข้อความ+สี), ใช้ @State เก็บค่าอินเด็กซ์ และเวลาแตะใช้ (index + 1) \% 3 ทำเสร็จส่งให้เพื่อนเล่นดู รับคำว้าวแรกว่า “โห ทำแอปเป็นแล้วเหรอ”

\end{tcbpractice}

\begin{tcbsummary}{สรุปบทนี้}
\begin{itemize}
  \item SwiftUI เป็นแบบประกาศ: คุณบรรยายว่า'หน้าตาเป็นอย่างไร' ระบบจัดการวาดและวาดใหม่ให้
  \item วิว = struct ที่คอนฟอร์มกับโปรโตคอล View; body ใช้บรรยายอินเทอร์เฟซ; modifier ตกแต่งแบบต่อกันเป็นโซ่
  \item วิวพื้นฐานสี่ตัว: Text / Button / TextField / Image; คลังไอคอน SF Symbols ฟรีเพียบ ใช้ไม่อั้น
  \item ชุดจัดวางสามพี่น้อง: VStack แนวตั้ง, HStack แนวนอน, ZStack ซ้อนกัน; Spacer คือพื้นที่ยืดหยุ่น
  \item @State ทำให้อินเทอร์เฟซมีชีวิต; \$ ไบนด์สองทางป้อนให้ช่องกรอก
  \item \#Preview เห็นอะไรได้อย่างนั้น; สีของระบบปรับเข้ากับโหมดมืดอัตโนมัติ; withAnimation ใส่ลูกเล่นได้ในบรรทัดเดียว
\end{itemize}

\end{tcbsummary}
